% pcra_future_work.tex
% Shelved content: Future Work - PCRA Package Conversion
% Removed from main proposal; kept here for future expansion.
%
% To compile standalone: pdflatex pcra_future_work.tex
%
\documentclass[11pt]{article}
\usepackage[margin=1in]{geometry}
\usepackage{booktabs,tabularx,amsmath,amssymb,listings,hyperref,url,xcolor,enumitem}
\definecolor{codegray}{rgb}{0.95,0.95,0.95}
\lstdefinestyle{python}{language=Python,basicstyle=\ttfamily\small,
  backgroundcolor=\color{codegray},breaklines=true,frame=single}
\title{Future Work: PCRA Package Conversion\\
\large Companion to the \texttt{robstatpy} GSoC 2026 Proposal}
\author{Aakarsh Gupta}
\date{\today}
\begin{document}
\maketitle

% SECTION 7 -- FUTURE WORK: PCRA
%--------------------------------------------------------------------
\section{Future Work: PCRA Package Conversion}

The \textbf{PCRA} (Portfolio Construction and Risk Analysis) R package,
also authored by Prof.\ Martin, depends on \texttt{RobStatTM} for
robust covariance inputs to portfolio optimisation and for robust
single-factor model fitting.
A Python equivalent of PCRA is a natural downstream application of
\texttt{robstatpy}.

\subsection{Dependency Relationships}

PCRA orchestrates 12+ R packages.  The critical dependency chain is:

\begin{center}
\begin{tabular}{ccccc}
  \fbox{PCRA} &
  $\xrightarrow{\text{covRobMM}}$ &
  \fbox{RobStatTM} &
  $\xrightarrow{\text{pyinit}}$ &
  \fbox{Fortran core} \\[6pt]
  $\downarrow$ & & & & \\
  \fbox{PortfolioAnalytics} &
  $\rightarrow$ &
  \fbox{PerformanceAnalytics} &
  $\rightarrow$ &
  \fbox{xts/zoo} \\
\end{tabular}
\end{center}

\noindent
For the Python side, the mapping is:

\renewcommand{\arraystretch}{1.2}
\begin{tabularx}{\linewidth}{@{} l l X @{}}
  \toprule
  \textbf{R Package} & \textbf{Python Equivalent} & \textbf{Notes} \\
  \midrule
  RobStatTM            & \texttt{robstatpy}        & This project \\
  PortfolioAnalytics   & \texttt{cvxpy}            & Convex optimisation \\
  PerformanceAnalytics & \texttt{empyrical}        & Risk/performance measures \\
  xts / zoo            & \texttt{pandas}           & Time series indexing \\
  ggplot2              & \texttt{matplotlib}        & Plotting \\
  \bottomrule
\end{tabularx}
\renewcommand{\arraystretch}{1}

\subsection{Preliminary Results}

A prototype Python reproduction of PCRA Chapter~2 has already been
completed (\texttt{pcra/python/ch2\_foundations\_demo.py}), demonstrating:

\begin{itemize}[leftmargin=1.5em, itemsep=2pt]
  \item Two-asset efficient frontier (Merton closed-form).
  \item Global minimum variance portfolio computation.
  \item Drawdown analysis and risk measures (VaR, ES).
  \item Robust location estimation via bisquare $\psi$-function.
\end{itemize}

\noindent
\textbf{Code example: efficient frontier computation (from the prototype):}

\begin{lstlisting}[style=python, caption={PCRA Chapter 2: Merton efficient frontier in Python}]
def math_efront_risky_mu_cov(mu_ret, vol_ret, corr_ret,
                              npoints=100):
    """Efficient frontier via closed-form Merton formula.
    Reproduces R's mathEfrontRiskyMuCov() from PCRA."""
    V = np.diag(vol_ret) @ corr_ret @ np.diag(vol_ret)
    Vinv = np.linalg.inv(V)
    one = np.ones(len(mu_ret))
    a = mu_ret @ Vinv @ one
    b = mu_ret @ Vinv @ mu_ret
    c = one @ Vinv @ one
    d = b * c - a ** 2
    mu_gmv = a / c
    sigma_gmv = 1 / np.sqrt(c)
    mu_range = np.linspace(mu_gmv, max(mu_ret)*1.2, npoints)
    sigma_range = np.sqrt(
        (c*mu_range**2 - 2*a*mu_range + b) / d)
    return mu_range, sigma_range, mu_gmv, sigma_gmv
\end{lstlisting}

% [Drawdown/VaR/ES code in pcra/python/ch2\_foundations\_demo.py]


\noindent
\textbf{Selected PCRA Chapter~2 results (Python vs.\ R):}

\renewcommand{\arraystretch}{1.2}
\begin{tabularx}{\linewidth}{@{} l r r c @{}}
  \toprule
  \textbf{Quantity} & \textbf{R} & \textbf{Python} & \textbf{Match} \\
  \midrule
  GMV $\sigma_P$ (2-asset, $\rho=0$)    & 0.1200 & 0.1200 & \checkmark \\
  GMV $\mu_P$ (2-asset, $\rho=0$)       & 0.0610 & 0.0610 & \checkmark \\
  3-asset GMV return                     & 0.0540 & 0.0540 & \checkmark \\
  3-asset GMV volatility                 & 0.0707 & 0.0707 & \checkmark \\
  VaR (5\%, SPY monthly)                 & matches & matches & \checkmark \\
  ES (5\%, SPY monthly)                  & matches & matches & \checkmark \\
  Max drawdown depth                     & matches & matches & \checkmark \\
  \bottomrule
\end{tabularx}
\renewcommand{\arraystretch}{1}

\vspace{0.3em}
\noindent
The complete comparison output (11~figures, 8~tables) is available in the
project repository:
\href{https://github.com/Aakarsh751/robstattm-pyport-AG-prop/blob/main/pcra/output/Ch2\_Comparison.pdf}{\texttt{pcra/output/Ch2\_Comparison.pdf}} (R vs.\ Python side-by-side).
These results match the R originals and validate the feasibility of
converting PCRA's core mathematical layer to Python.
Full PCRA conversion is scoped as a stretch goal; the foundation
laid by \texttt{robstatpy} makes it achievable in a follow-up
contribution.


\noindent
The 	\texttt{rpy2} wrapper methodology scales directly to PCRA and other
R packages in the robust statistics ecosystem: once conversion patterns
are established for 	\texttt{RobStatTM}, extending them to PCRA is
incremental. Full PCRA conversion is scoped as a stretch goal, contingent
on completing the core 	\texttt{robstatpy} deliverables.


\end{document}
